% !TeX root = ../mathematical-methods1.tex
\chapter{Partial Derivatives}\label{chap:partial-derivatives}

This chapter concerns derivatives of functions of two or more variables,
also called multivariable derivatives. It begins on handwritten p.~69 of
\texttt{Slides\_Week-5.pdf}. Functions are assumed to be ``nice enough'',
meaning that the derivatives used exist and are continuous.

\section{Introduction}

For $f:(a,b)\to\mathbb R$, the graph of $f$ is a \emph{curve} $y=f(x)$ in the
$(x,y)$-plane. For a function of two variables $f(x,y)$, the graph
$z=f(x,y)$ is a \emph{surface} in $(x,y,z)$-space.

\begin{definition}[Level sets]
For $c\in\mathbb R$, the level set of $f$ at height $c$ is
\[
 \{(x,y): f(x,y)=c\}.
\]
Drawing several level sets in the $(x,y)$-plane gives a contour map of the surface.
\end{definition}

\begin{example}[Lecture Example 1: a paraboloid]
Sketch the surface $z=f(x,y)$ for $f(x,y)=x^2+y^2$.

The level sets are:
\begin{itemize}
 \item $c<0$: the empty set $\varnothing$;
 \item $c=0$: the single point $x=y=0$;
 \item $c>0$: $x^2+y^2=c$, a circle of radius $\sqrt c$ centred at the origin.
\end{itemize}
Lifting each circle to height $z=c$ (an ``elevated map'') produces the
paraboloid $z=x^2+y^2$.
\end{example}

\begin{center}
\begin{tikzpicture}[>=Stealth]
 \begin{scope}
  \draw[->] (-2,0)--(2.1,0) node[right] {$x$};
  \draw[->] (0,-2)--(0,2.1) node[above] {$y$};
  \fill[ink] (0,0) circle (1.5pt);
  \foreach \c/\r in {1/0.7,2/0.99,3/1.21,4/1.4}
   \draw[ink] (0,0) circle (\r);
  \node[font=\small] at (0,-2.5) {Level sets $x^2+y^2=c$};
 \end{scope}
 \begin{scope}[xshift=6cm,yshift=-1.6cm]
  \draw[->] (0,0)--(0,3.5) node[above] {$z$};
  \draw[->] (0,0)--(2.2,0) node[right] {$y$};
  \draw[->] (0,0)--(-1.2,-0.9) node[below left] {$x$};
  \draw[ink,thick,domain=-1.6:1.6,samples=50] plot (\x,{\x*\x});
  \foreach \h in {0.64,1.44,2.56}
   \draw[blue] (0,\h) ellipse ({sqrt(\h)} and {0.18*sqrt(\h)});
  \node[font=\small] at (0.4,-1.5) {Paraboloid $z=x^2+y^2$};
 \end{scope}
\end{tikzpicture}
\end{center}

For $y=f(x)$, the tangent line to the graph at $x_0$ is
$L_t:\ y=f(x_0)+f'(x_0)(x-x_0)$. For a surface $z=f(x,y)$, the analogue is a
\emph{tangent plane} $P_t$ at $(x_0,y_0)$. To describe $P_t$ we need two
vectors in it, which we obtain from \emph{sections} (or \emph{slices}) of the
surface:
\begin{itemize}
 \item the $y$-slice ($y=y_0$ held constant): $z=f(x,y_0)$, a function of $x$;
 \item the $x$-slice ($x=x_0$ held constant): $z=f(x_0,y)$, a function of $y$.
\end{itemize}
The slopes of the tangent lines to these slices are the partial derivatives.

\section{Partial derivatives}

\begin{definition}[Partial derivatives]
The partial derivatives of $f(x,y)$ at $(x_0,y_0)$ are
\begin{align*}
 \frac{\partial f}{\partial x}(x_0,y_0)
 &:=\lim_{h\to0}\frac{f(x_0+h,y_0)-f(x_0,y_0)}{h},\\
 \frac{\partial f}{\partial y}(x_0,y_0)
 &:=\lim_{h\to0}\frac{f(x_0,y_0+h)-f(x_0,y_0)}{h},
\end{align*}
provided the limits exist.
\end{definition}

The symbol $\partial$ is read ``curly d''. As $f'$ abbreviates
$\mathrm df/\mathrm dx$, we abbreviate
\[
 \frac{\partial f}{\partial x}=f_x,\qquad \frac{\partial f}{\partial y}=f_y.
\]
To calculate a partial derivative, use the usual single-variable rules with
respect to the variable in question, \emph{holding the other variable(s) constant}.

\begin{example}[Lecture Example 2]
Let $f(x,y)=\sin(x^2y)+xy$. Find $f_x$ and $f_y$.

Holding $y$ constant and using the chain rule,
\[
 f_x=\cos(x^2y)\cdot2xy+y.
\]
Holding $x$ constant,
\[
 f_y=\cos(x^2y)\cdot x^2+x.
\]
\end{example}

\subsection{Higher derivatives}
The second derivatives are
\[
 \frac{\partial^2f}{\partial x^2}=f_{xx}:=\frac{\partial}{\partial x}\Bigl(\frac{\partial f}{\partial x}\Bigr)=(f_x)_x,
 \qquad
 \frac{\partial^2f}{\partial y^2}=f_{yy}:=\frac{\partial}{\partial y}\Bigl(\frac{\partial f}{\partial y}\Bigr)=(f_y)_y,
\]
together with the \emph{mixed derivatives}
\[
 \frac{\partial^2f}{\partial y\,\partial x}:=\frac{\partial}{\partial y}\Bigl(\frac{\partial f}{\partial x}\Bigr)=(f_x)_y=f_{xy},
 \qquad
 f_{yx}:=(f_y)_x.
\]
In the subscript notation, differentiate in the order the subscripts are read
from left to right.

\begin{example}[Lecture Example 3]
Find all first and second derivatives of $f(x,y)=\sin(x^2y)+xy$.

From Example~2,
\[
 f_x=2xy\cos(x^2y)+y,\qquad f_y=x^2\cos(x^2y)+x.
\]
The second derivatives are
\begin{align*}
 f_{xx}&=2y\cos(x^2y)-4x^2y^2\sin(x^2y),\\
 f_{yy}&=-x^4\sin(x^2y),\\
 f_{xy}&=2x\cos(x^2y)-2x^3y\sin(x^2y)+1,\\
 f_{yx}&=2x\cos(x^2y)-2x^3y\sin(x^2y)+1.
\end{align*}
\end{example}

Observe that $f_{xy}=f_{yx}$. This holds generally.

\begin{theorem}[Clairaut's theorem]
If $f_x$, $f_y$, $f_{xx}$, $f_{xy}$, $f_{yx}$ and $f_{yy}$ exist and are
continuous, then
\[
 f_{xy}=f_{yx}.
\]
\end{theorem}
The proof belongs to Analysis. More generally, for smooth enough $f$ the
order of differentiation does not matter for any higher derivatives, for example
$f_{xxy}=f_{xyx}=f_{yxx}$.

\section{Tangent plane and linear approximation}

For $y=f(x)$, the slope $f'(x_0)$ gives a direction vector
$\mathbf t=\mathbf i+f'(x_0)\mathbf j$ for the tangent line.
For $z=f(x,y)$, the partial derivatives $f_x$ and $f_y$ are the slopes of the
tangent vectors to the $y$- and $x$-slices.

The $y$-slice at $y=y_0$ is the curve $\mathbf r(x)=\langle x,\,y_0,\,f(x,y_0)\rangle$.
Its tangent vector
\[
 \mathbf u=\langle1,\,0,\,f_x(x_0,y_0)\rangle
\]
lies in the tangent plane $P_t$. Similarly, the tangent to the $x$-slice,
\[
 \mathbf v=\langle0,\,1,\,f_y(x_0,y_0)\rangle,
\]
lies in $P_t$.

\begin{center}
\begin{tikzpicture}[>=Stealth]
 % axes: z up, y right, x towards the viewer (down-left)
 \draw[->] (0,0)--(0,3.6) node[above] {$z$};
 \draw[->] (0,0)--(4.6,0) node[right] {$y$};
 \draw[->] (0,0)--(-1.8,-1.5) node[below left] {$x$};
 % tangent plane as a parallelogram
 \draw[muted,fill=ink!6] (0.2,1.3)--(3.6,1.9)--(2.4,3.4)--(-1.0,2.8)--cycle;
 \node[font=\small] at (3.4,3.1) {$P_t$};
 % point of tangency and its foot (x0,y0)
 \coordinate (P) at (1.3,2.35);
 \coordinate (F) at (1.3,-0.6);
 \draw[dashed,muted] (F)--(P);
 \draw[dashed,muted] (F)--(-0.72,-0.6) node[left] {$x_0$};
 \draw[dashed,muted] (F)--(2.02,0) node[above right] {$y_0$};
 \fill (P) circle (1.6pt);
 \draw[->,thick,red] (P)--++(-0.75,-0.45) node[left] {$\mathbf u$};
 \draw[->,thick,blue] (P)--++(1.1,0.2) node[right] {$\mathbf v$};
\end{tikzpicture}
\end{center}

A parametric equation of $P_t$ is therefore
\[
 \mathbf r=\mathbf r_0+s\mathbf u+t\mathbf v,\qquad s,t\in\mathbb R,
 \qquad \mathbf r_0=\langle x_0,\,y_0,\,f(x_0,y_0)\rangle.
\]
A normal to $P_t$ is
\[
 \mathbf n=\mathbf u\times\mathbf v=
 \begin{vmatrix}
  \mathbf i&\mathbf j&\mathbf k\\
  1&0&f_x\\
  0&1&f_y
 \end{vmatrix}
 =\langle-f_x,\,-f_y,\,1\rangle.
\]
A Cartesian equation of $P_t$ is $\mathbf n\cdot(\mathbf r-\mathbf r_0)=0$, that is,
\[
 -f_x(x-x_0)-f_y(y-y_0)+\bigl(z-f(x_0,y_0)\bigr)=0.
\]
Rearranging,
\begin{equation}
 \boxed{z=f(x_0,y_0)+f_x(x_0,y_0)(x-x_0)+f_y(x_0,y_0)(y-y_0)}.
 \tag{1}\label{eq:mm-tangent-plane}
\end{equation}
Equation \eqref{eq:mm-tangent-plane} is the equation of the tangent plane
$P_t$ at $(x_0,y_0)$. It is also the first-order approximation of $f(x,y)$
near $(x_0,y_0)$:
\[
 f(x,y)\approx f(x_0,y_0)+f_x(x_0,y_0)(x-x_0)+f_y(x_0,y_0)(y-y_0).
\]

\begin{example}[Lecture Example 4]
Find the tangent plane to $z=f(x,y)=2x^2+y^2$ at $(x_0,y_0)=(1,1)$.

Here $f(1,1)=2+1=3$, $f_x=4x$ so $f_x(1,1)=4$, and $f_y=2y$ so $f_y(1,1)=2$.
By \eqref{eq:mm-tangent-plane},
\[
 z=3+4(x-1)+2(y-1)
 \quad\Longleftrightarrow\quad
 \boxed{z=4x+2y-3}.
\]
\end{example}

\section{The multivariable chain rule}

Recall the chain rule for $y=y(x)$ and $x=x(t)$:
\[
 \frac{\mathrm d}{\mathrm dt}y(x(t))=y'(x(t))\,x'(t)
 \quad\Longleftrightarrow\quad
 \frac{\mathrm dy}{\mathrm dt}=\frac{\mathrm dy}{\mathrm dx}\frac{\mathrm dx}{\mathrm dt}.
\]
In two variables, let $z=f(x,y)$ with $x=x(t)$ and $y=y(t)$.
What is $\frac{\mathrm d}{\mathrm dt}f(x(t),y(t))$? It follows from Analysis that
\[
 \frac{\mathrm d}{\mathrm dt}f(x(t),y(t))
 =f_x\bigl(x(t),y(t)\bigr)\,x'(t)+f_y\bigl(x(t),y(t)\bigr)\,y'(t),
\]
or, in Leibniz notation,
\begin{equation}
 \boxed{\frac{\mathrm d}{\mathrm dt}f(x(t),y(t))
 =\frac{\partial f}{\partial x}\frac{\mathrm dx}{\mathrm dt}
 +\frac{\partial f}{\partial y}\frac{\mathrm dy}{\mathrm dt}}.
 \tag{2}\label{eq:mm-chain}
\end{equation}

\textbf{Exercise.} Derive \eqref{eq:mm-chain} formally using the linear
approximations of $f(x,y)$ and of $x(t)$, $y(t)$.

\begin{example}[Lecture Example 5]
Let $z=f(x,y)=x^2y+3xy^4$ with $x(t)=\sin(2t)$ and $y(t)=\cos t$.
Find $\mathrm df/\mathrm dt$ in terms of $t$.

We have
\[
 f_x=2xy+3y^4,\quad f_y=x^2+12xy^3,\quad x'(t)=2\cos2t,\quad y'(t)=-\sin t.
\]
By \eqref{eq:mm-chain},
\begin{align*}
 \frac{\mathrm df}{\mathrm dt}
 &=(2xy+3y^4)\,2\cos2t+(x^2+12xy^3)(-\sin t)\\
 &=2\cos2t\bigl(2\sin2t\cos t+3\cos^4t\bigr)
   -\sin t\bigl(\sin^22t+12\sin2t\cos^3t\bigr).
\end{align*}
\end{example}

\subsection{More general chain rules}
Now let $z=f(x,y)$ with $x=x(u,v)$ and $y=y(u,v)$. This is a multivariable
change of variables $(x,y)\leftrightarrow(u,v)$. We want
\[
 f_u:=\frac{\partial}{\partial u}f\bigl(x(u,v),y(u,v)\bigr),
 \qquad
 f_v:=\frac{\partial}{\partial v}f\bigl(x(u,v),y(u,v)\bigr).
\]
Applying \eqref{eq:mm-chain} with $v$ held fixed for $f_u$, and with $u$ held
fixed for $f_v$, gives
\begin{equation}
 \boxed{
 \frac{\partial z}{\partial u}=\frac{\partial z}{\partial x}\frac{\partial x}{\partial u}
 +\frac{\partial z}{\partial y}\frac{\partial y}{\partial u},
 \qquad
 \frac{\partial z}{\partial v}=\frac{\partial z}{\partial x}\frac{\partial x}{\partial v}
 +\frac{\partial z}{\partial y}\frac{\partial y}{\partial v}}.
 \tag{3}\label{eq:mm-chain-uv}
\end{equation}
In matrix form, \eqref{eq:mm-chain-uv} reads
\begin{equation}
 \begin{pmatrix} z_u\\ z_v\end{pmatrix}
 =\begin{pmatrix} x_u&y_u\\ x_v&y_v\end{pmatrix}
  \begin{pmatrix} z_x\\ z_y\end{pmatrix}.
 \tag{3'}
\end{equation}
Here
\[
 \begin{pmatrix} x_u&y_u\\ x_v&y_v\end{pmatrix}
 =\begin{pmatrix}
  \dfrac{\partial x}{\partial u}&\dfrac{\partial y}{\partial u}\\[8pt]
  \dfrac{\partial x}{\partial v}&\dfrac{\partial y}{\partial v}
 \end{pmatrix}
\]
is the \emph{Jacobian matrix} of the change of variables, and
\[
 \nabla f=\begin{pmatrix} f_x\\ f_y\end{pmatrix}
 =\begin{pmatrix} \partial f/\partial x\\ \partial f/\partial y\end{pmatrix}
\]
is the \emph{gradient vector} of $f$, or ``gradient of $f$'' for short.
It is also written $\boldsymbol\nabla f$. The gradient has many uses.

\begin{remark}
Some texts write the Jacobian as the transpose of this matrix, with rows
indexed by $x,y$ and columns by $u,v$. The lecture's arrangement makes
\eqref{eq:mm-chain-uv} a matrix times the gradient.
\end{remark}

\section{Directional derivatives and the gradient}

Recall that $f_x$ and $f_y$ are the slopes of $f$ in the $x$- and
$y$-directions, along the slices. What is the slope in some other direction?

\begin{definition}[Directional derivative]
Let $\hat{\mathbf u}=\langle u_1,u_2\rangle$ be a unit vector in the $(x,y)$-plane.
The directional derivative of $f(x,y)$ in the direction $\hat{\mathbf u}$ is
the slope of the slice in that direction:
\[
 D_{\hat{\mathbf u}}f(x_0,y_0)
 :=\lim_{h\to0}\frac{f(x_0+hu_1,\,y_0+hu_2)-f(x_0,y_0)}{h}.
\]
\end{definition}

Equivalently,
\[
 D_{\hat{\mathbf u}}f(x_0,y_0)
 =\frac{\mathrm d}{\mathrm dt}\Bigl[f(x_0+tu_1,\,y_0+tu_2)\Bigr]_{t=0}.
\]
By the chain rule \eqref{eq:mm-chain}, with $x'(t)=u_1$ and $y'(t)=u_2$,
\[
 D_{\hat{\mathbf u}}f(x_0,y_0)
 =f_x(x_0,y_0)\,u_1+f_y(x_0,y_0)\,u_2
 =\langle f_x,f_y\rangle\cdot\langle u_1,u_2\rangle.
\]
Hence
\[
 \boxed{D_{\hat{\mathbf u}}f(x_0,y_0)=\nabla f(x_0,y_0)\cdot\hat{\mathbf u}},
 \qquad
 \nabla f:=\begin{pmatrix} \partial f/\partial x\\ \partial f/\partial y\end{pmatrix}.
\]
Taking $\hat{\mathbf u}=\mathbf i$ or $\mathbf j$ recovers $f_x$ or $f_y$.

\subsection{Interpreting the gradient}
The value $|D_{\hat{\mathbf u}}f(x_0,y_0)|$ is the magnitude of the slope in
the direction $\hat{\mathbf u}$, and its sign tells us:
\begin{itemize}
 \item $D_{\hat{\mathbf u}}f>0$: $f$ increases in the direction $\hat{\mathbf u}$;
 \item $D_{\hat{\mathbf u}}f<0$: $f$ decreases in the direction $\hat{\mathbf u}$;
 \item $D_{\hat{\mathbf u}}f=0$: $f$ is approximately constant, so
       $\hat{\mathbf u}$ is a level-set direction.
\end{itemize}
Let $\theta$ be the angle between $\nabla f$ and $\hat{\mathbf u}$. Since
$|\hat{\mathbf u}|=1$,
\[
 D_{\hat{\mathbf u}}f=\nabla f\cdot\hat{\mathbf u}=|\nabla f|\cos\theta,
 \qquad |D_{\hat{\mathbf u}}f|\le|\nabla f|.
\]
Therefore, when $\nabla f\ne\mathbf 0$:
\begin{itemize}
 \item $\theta=0$: $\hat{\mathbf u}$ is the direction of \emph{steepest ascent},
       with slope $|\nabla f|$;
 \item $\theta=\pi$: $\hat{\mathbf u}$ is the direction of \emph{steepest descent},
       with slope $-|\nabla f|$;
 \item $\theta=\pi/2$: $\hat{\mathbf u}\perp\nabla f$, a level-set direction.
\end{itemize}

\begin{center}
\fbox{\parbox{0.8\linewidth}{\centering
$\nabla f$ points in the direction of steepest ascent, and is perpendicular to the level sets of $f$.}}
\end{center}

\begin{center}
\begin{tikzpicture}[>=Stealth]
 \foreach \r in {0.7,1.2,1.7,2.2}
  \draw[ink!60] (0,0) circle (\r);
 \fill (1.2,0) circle (1.5pt);
 \draw[->,thick,red] (1.2,0)--(2.2,0) node[right] {$\nabla f$};
 \draw[->,thick,blue] (1.2,0)--(1.2,1) node[above right,fill=white,inner sep=1pt,font=\small] {level-set direction};
 \node[font=\small] at (0,-2.7) {Level sets of $x^2+y^2$ with the gradient at a point};
\end{tikzpicture}
\end{center}
